\chapter*{Annexes}
\addcontentsline{toc}{chapter}{Annexes}
\markboth{ANNEXES}{ANNEXES}

\section*{Annexe A : Catalogue des concepts ontologiques du Knowledge Graph}

Le tableau \ref{tab:annexe_concepts} présente un extrait représentatif de l'ontologie des 39 concepts interconnectés structurés dans le module \texttt{concept\_catalog.py} de notre studio.

\begin{table}[H]
\centering
\caption{Extrait de l'ontologie des concepts du Knowledge Graph}
\label{tab:annexe_concepts}
\begin{tabularx}{\textwidth}{|l|l|l|X|}
\hline
\rowcolor{focpgreen!18}
\textbf{Clé du Concept} & \textbf{Domaine} & \textbf{Importance} & \textbf{Description Métier} \\
\hline
\texttt{business.objective} & business & Critique & Objectif métier central et raison d'être du système. \\
\hline
\texttt{business.success\_metrics} & business & Haute & Indicateurs clés de performance et cibles chiffrées. \\
\hline
\texttt{users.primary\_personas} & users & Critique & Profils utilisateurs principaux et segmentation des accès. \\
\hline
\texttt{beneficiaries.categories} & beneficiaries & Critique & Répartition des bénéficiaires (femmes, jeunes, handicap). \\
\hline
\texttt{entities.core\_objects} & entities & Critique & Entités de gestion centrales (coopératives, conventions). \\
\hline
\texttt{security.auth\_method} & security & Critique & Mécanisme d'authentification et gestion des sessions. \\
\hline
\texttt{security.data\_class} & security & Haute & Classification des données (sensibles, publiques, restreintes). \\
\hline
\texttt{deployment.target\_env} & deployment & Haute & Environnement cible et topologie d'infrastructure. \\
\hline
\texttt{monitoring.kpis} & monitoring & Moyenne & Métriques d'exploitation et tableaux de bord décisionnels. \\
\hline
\end{tabularx}
\end{table}

\section*{Annexe B : Règles d'audit du Moteur de Revue Qualité}

Le Moteur de Revue Qualité évalue plus de 100 règles dérivées de 90 gabarits de vérification, résumés dans le tableau \ref{tab:annexe_rules}.

\begin{table}[H]
\centering
\caption{Catégories des règles de validation du Moteur d'Audit}
\label{tab:annexe_rules}
\begin{tabularx}{\textwidth}{|l|c|X|}
\hline
\rowcolor{focpgreen!18}
\textbf{Famille de Règles} & \textbf{Nombre} & \textbf{Critères de Validation et d'Homologation} \\
\hline
Complétude conceptuelle & 39 & Vérifie que chaque concept a atteint un seuil de complétude $\ge 75\%$. \\
\hline
Confiance ontologique & 39 & Évalue que l'indice de confiance des données capturées est $\ge 70\%$. \\
\hline
Couverture des domaines & 15 & S'assure de la complétude minimale des 15 domaines clés du projet. \\
\hline
Présence des types d'exigences & 11 & Contrôle l'existence d'exigences BRD, FRD, NFR et sécurité. \\
\hline
Contrôles de cybersécurité & 8 & Vérifie l'application effective des contrôles essentiels (RBAC, TLS, journalisation, chiffrement). \\
\hline
Cohérence structurelle & 11 & Valide les relations croisées, l'absence de noeuds orphelins et les critères d'acceptation. \\
\hline
\end{tabularx}
\end{table}

\section*{Annexe C : Extrait de configuration Docker Compose}

\begin{lstlisting}[language=bash, caption={Configuration Docker Compose multi-services}, label={lst:docker_compose}]
name: srs-platform

services:
  db:
    image: postgres:16-alpine
    restart: unless-stopped
    environment:
      POSTGRES_USER: srs
      POSTGRES_PASSWORD: srs
      POSTGRES_DB: srs
    ports:
      - "5432:5432"

  backend:
    build:
      context: ./backend
      dockerfile: Dockerfile
    environment:
      DATABASE_URL: postgresql+asyncpg://srs:srs@db:5432/srs
      REDIS_URL: redis://redis:6379/0
      AI_PROVIDER: gemini
    depends_on:
      - db
    ports:
      - "8000:8000"

  frontend:
    build:
      context: ./frontend
      dockerfile: Dockerfile
    environment:
      NEXT_PUBLIC_API_BASE_URL: http://localhost:8000/api/v1
    depends_on:
      - backend
    ports:
      - "3000:3000"
\end{lstlisting}

\newpage
